% ============================================
% Johns Hopkins Letter Template
% Assistant/Associate Professor Cover Letter – University of Texas at Austin
% School of Civic Leadership
% ============================================

\documentclass[11pt,letterpaper]{letter}

\usepackage{graphicx}
\usepackage{eso-pic}
\usepackage{geometry}
\usepackage{microtype}
\usepackage[hidelinks]{hyperref}

% ----- Page Setup -----
\geometry{left=25mm,right=25mm,top=25mm,bottom=25mm}

% ----- Logo Placement (foreground, first page only) -----
\newcommand{\PutJHULogoFG}[1]{%
  \AddToShipoutPictureFG*{%
    \AtPageUpperLeft{%
      \hspace{10mm}%
      \raisebox{-38mm}{%
        \includegraphics[width=6cm]{#1}%
      }%
    }%
  }%
}

% ----- Sender Information -----
\signature{Decory Edwards}
\address{Department of Economics \\ Johns Hopkins University \\ Baltimore, MD 21218 \\ \texttt{dedwar65@jh.edu}}

\begin{document}
\PutJHULogoFG{Images/jhulogo.png}

\begin{letter}{Search Committee \\
School of Civic Leadership \\
The University of Texas at Austin \\
Austin, TX 78712 \\ USA}

\opening{Dear Members of the Search Committee,}

I am writing to express my interest in the open-rank faculty position in Economics at the School of Civic Leadership at the University of Texas at Austin, beginning Fall 2026. I am a Ph.D. candidate in Economics at Johns Hopkins University, advised by Professors Christopher D. Carroll, Nicholas W. Papageorge, and Stelios Fourakis, and I expect to receive my degree in May 2026. My research fields are macroeconomics, wealth and income dynamics, and computational economics.

My research agenda focuses on understanding the determinants of wealth inequality and the mechanisms through which heterogeneity in returns to wealth shapes aggregate outcomes. In my job-market paper, I estimate the degree of heterogeneity in individual portfolio returns using micro-data from the Survey of Consumer Finances and simulate a structural model in which households face idiosyncratic return risk. I show that allowing for persistent heterogeneity in returns substantially improves the model’s ability to match the empirical distribution of wealth, offering a tractable way to connect micro-level return dispersion to macroeconomic inequality.

In a second project, I use the Health and Retirement Study to examine the relationship between interpersonal trust and portfolio performance. Building on the literature that finds a hump-shaped relationship between trust and income, I show that a similar relationship holds for individual returns to wealth measured in the HRS data, highlighting a behavioral channel through which social attitudes shape saving and investment outcomes. This work also provides new estimates of the persistent component of returns to net worth, which motivate the calibration of heterogeneity in my structural model.

The mission of the School of Civic Leadership to deepen understanding of how markets, institutions, and incentives sustain a free and prosperous society resonates directly with my research. My work contributes to this broader inquiry by studying how individual decision-making under uncertainty aggregates into macroeconomic outcomes and shapes economic opportunity across households. I am especially drawn to the School’s interdisciplinary orientation and to the Civitas Institute’s role in fostering research that connects economic reasoning to civic responsibility and institutional design. My background in macroeconomics and household finance naturally complements such a dialogue between economics, public policy, and civic thought.

\textbf{Teaching.} My teaching experience centers on undergraduate macroeconomics and an upper-level macro-policy seminar, where I have led weekly discussion sections and held regular office hours for classes ranging from 16 to 40 students. My approach is student-centered: I aim to help students “convince themselves” that they have moved from confusion to understanding by holding structured office-hour coaching that builds trust and confidence. At UT Austin, I am prepared to teach \textit{Principles of Macroeconomics}, \textit{Intermediate Macroeconomic Theory}, \textit{Money and Banking}, and \textit{Quantitative Methods for Economics}. I would also welcome developing electives such as \textit{Economics of Inequality and Tax Policy} or \textit{Computational Macroeconomics and Policy}, emphasizing reproducible workflows and evidence-based decision-making.

I view teaching and research as complementary ways of preparing students for leadership in a free society. In the classroom, I aim to demonstrate how rigorous economic reasoning can illuminate the institutional foundations of prosperity and civic life. In research, I seek to build empirically grounded models that connect individual behavior and aggregate outcomes in ways that inform policy and broaden participation in economic opportunity.

I would be honored to contribute to the School of Civic Leadership’s dual mission of excellence in scholarship and civic education, and to collaborate with colleagues and students in advancing the understanding of how economic institutions and incentives shape a free and prosperous society. Thank you for your consideration; I welcome the opportunity to discuss my fit with your School at your convenience.

\closing{Sincerely,}

\end{letter}
\end{document}
